\chapter{Contraction Mappings}
\label{ch:ii16}

The contraction mapping theorem turns a quantitative shrinking estimate into existence, uniqueness, and an algorithm. It is one of the cleanest bridges between analysis and computation.

\section*{Learning goals}
After this chapter, the reader should be able to:
\begin{itemize}
\item verify that a self-map is a contraction by finding a constant \(q<1\);
\item prove existence and uniqueness of a fixed point by applying the contraction mapping theorem with all hypotheses checked;
\item derive quantitative error bounds for Picard iteration;
\item choose a closed invariant subset on which a proposed iteration is actually contractive;
\item compare different fixed-point iterations by their contraction constants;
\item construct examples showing why completeness or the strict contraction condition cannot simply be omitted;
\end{itemize}
\section*{Conceptual roadmap}
\[
\boxed{\text{definition}\;\longrightarrow\;\text{estimate}\;\longrightarrow\;\text{theorem}\;\longrightarrow\;\text{application}.}
\]

\section{Contractions}
A contraction reduces all distances by a fixed factor $q<1$.

\section{Fixed points}
A fixed point solves $Tx=x$.

\section{Picard iteration}
Repeated application $x_{n+1}=Tx_n$ is the natural algorithm.


% BEGIN VOL02-EXPANSION II16-example-01
\begin{example}[Picard iteration with an exact fixed point check]\label{ex:ii16-expand-01}
Let
\[
T(x)=\frac{x+3}{4}.
\]
Then
\[
|T(x)-T(y)|=\frac14|x-y|,
\]
so \(q=1/4\). Starting from \(x_0=0\),
\[
x_1=\frac34,\qquad x_2=\frac{15}{16},\qquad x_3=\frac{63}{64}.
\]
The fixed-point equation gives \(x_*=1\), and indeed
\[
|x_n-1|=4^{-n}.
\]
The explicit formula checks the geometric convergence predicted by Banach's
theorem.
\end{example}
% END VOL02-EXPANSION II16-example-01
\section{Completeness}
The iteration is Cauchy because geometric distance estimates telescope; completeness supplies the limit.

\section{Uniqueness}
Two fixed points would be strictly closer after applying $T$, impossible unless equal.

\section{Error estimates}
A priori and a posteriori geometric bounds quantify iteration error.


% BEGIN VOL02-EXPANSION II16-example-02
\begin{example}[A posteriori error from the last step]\label{ex:ii16-expand-02}
For a contraction with factor \(q<1\), Picard iterates satisfy
\[
d(x_n,x_*)\le \frac{q}{1-q}d(x_n,x_{n-1}).
\]
Thus if \(q=0.2\) and the most recent step has size \(10^{-5}\), then
\[
d(x_n,x_*)\le\frac{0.2}{0.8}10^{-5}=2.5\times10^{-6}.
\]
This is computationally useful because it estimates the unknown fixed-point
error from quantities already produced by the iteration.
\end{example}
% END VOL02-EXPANSION II16-example-02
\section{Parameter dependence}
Contractive fixed points vary stably under perturbations.


% BEGIN VOL02-EXPANSION II16-example-03
\begin{example}[Stable dependence of a fixed point on a parameter]\label{ex:ii16-expand-03}
Consider
\[
T_\lambda(x)=\frac{x+\lambda}{3}.
\]
Every \(T_\lambda\) has contraction factor \(1/3\) and fixed point
\(x_\lambda=\lambda/2\). If \(\lambda,\mu\) are two parameters,
\[
|x_\lambda-x_\mu|=\frac12|\lambda-\mu|.
\]
The general perturbation theorem gives the compatible bound
\[
|x_\lambda-x_\mu|
\le\frac{1}{1-1/3}\sup_x|T_\lambda(x)-T_\mu(x)|
=\frac12|\lambda-\mu|.
\]
Here the stability estimate is sharp.
\end{example}
% END VOL02-EXPANSION II16-example-03
\section{Applications}
Integral equations and ODEs can be reformulated as fixed-point problems.

\section{Core structural results}

\begin{theorem}[Banach fixed-point theorem]\label{thm:ii16-01}
A contraction on a nonempty complete metric space has a unique fixed point, and every Picard iteration converges to it.
\begin{proof}
The successive differences decay geometrically, making iterates Cauchy. Completeness gives a limit; continuity of the contraction makes it fixed; the contraction inequality gives uniqueness.
\end{proof}
\end{theorem}

\begin{theorem}[A priori error bound]\label{thm:ii16-02}
If $x_{n+1}=Tx_n$ and contraction factor is $q$, then $d(x_n,x_*)\le q^n(1-q)^{-1}d(x_1,x_0)$.
\begin{proof}
Bound the tail by the geometric sum of successive differences.
\end{proof}
\end{theorem}

\begin{theorem}[Stability under perturbation]\label{thm:ii16-03}
If contractions $T,S$ share factor $q<1$ and fixed points $x_T,x_S$, then $d(x_T,x_S)\le(1-q)^{-1}\sup_x d(Tx,Sx)$.
\begin{proof}
Insert and subtract $Tx_S$ and use the contraction bound before rearranging.
\end{proof}
\end{theorem}

\section{Worked examples}

\begin{example}[Affine contraction]\label{ex:ii16-worked-01}
Solve $x=(x+3)/4$. The map has factor $1/4$ and fixed point $x=1$.
\end{example}

\begin{example}[Iteration rate]\label{ex:ii16-worked-02}
How many powers of $q$ control the error? Error decays geometrically like $q^n$, up to the prefactor determined by the first step.
\end{example}

\begin{example}[Uniqueness proof]\label{ex:ii16-worked-03}
Show two fixed points must coincide. $d(x,y)=d(Tx,Ty)\le qd(x,y)$; since $q<1$, the distance is zero.
\end{example}

\begin{example}[Why completeness matters]\label{ex:ii16-worked-04}
Give a contraction on an incomplete space with no fixed point. On $(0,1)$, $T(x)=(x+1)/2$ is a contraction whose fixed point $1$ lies outside the space.
\end{example}

\section{Solved dossiers}
The dossiers are longer solved problems. Legacy mappings guide topic selection where explicit retained source blocks exist; otherwise fresh canonical problems complete the chapter.

\begin{problem}[Affine contraction]\label{prob:ii16-dossier-01}
Solve $x=(x+3)/4$.
\end{problem}
\begin{solution}
The map has factor $1/4$ and fixed point $x=1$.
\end{solution}

\begin{problem}[Iteration rate]\label{prob:ii16-dossier-02}
How many powers of $q$ control the error?
\end{problem}
\begin{solution}
Error decays geometrically like $q^n$, up to the prefactor determined by the first step.
\end{solution}

\begin{problem}[Uniqueness proof]\label{prob:ii16-dossier-03}
Show two fixed points must coincide.
\end{problem}
\begin{solution}
$d(x,y)=d(Tx,Ty)\le qd(x,y)$; since $q<1$, the distance is zero.
\end{solution}

\begin{problem}[Why completeness matters]\label{prob:ii16-dossier-04}
Give a contraction on an incomplete space with no fixed point.
\end{problem}
\begin{solution}
On $(0,1)$, $T(x)=(x+1)/2$ is a contraction whose fixed point $1$ lies outside the space.
\end{solution}

\begin{problem}[Matrix contraction]\label{prob:ii16-dossier-05}
When is $T(x)=Ax+b$ a Euclidean contraction?
\end{problem}
\begin{solution}
Whenever the operator norm $\|A\|_2<1$.
\end{solution}

\begin{problem}[Fixed-point series]\label{prob:ii16-dossier-06}
Solve $x=Ax+b$ with $\|A\|<1$.
\end{problem}
\begin{solution}
$x=(I-A)^{-1}b=\sum_{n\ge0}A^n b$, with convergence from the geometric operator-norm bound.
\end{solution}

\begin{problem}[A posteriori estimate]\label{prob:ii16-dossier-07}
Bound $d(x_n,x_*)$ using $d(x_n,x_{n-1})$.
\end{problem}
\begin{solution}
The tail is at most $d(x_n,x_{n-1})/(1-q)$.
\end{solution}

\begin{problem}[Nested mapping]\label{prob:ii16-dossier-08}
Why must $T$ map the complete set into itself?
\end{problem}
\begin{solution}
Otherwise the iteration may leave the domain and the completeness argument no longer applies.
\end{solution}

\begin{problem}[Parameter stability]\label{prob:ii16-dossier-09}
What happens if $T$ is perturbed by at most epsilon uniformly?
\end{problem}
\begin{solution}
The fixed point moves by at most $\varepsilon/(1-q)$.
\end{solution}

\begin{problem}[Integral-equation preview]\label{prob:ii16-dossier-10}
How does an integral operator become a contraction?
\end{problem}
\begin{solution}
Bound the sup-norm difference of outputs by an interval-length times a Lipschitz constant times the input sup difference.
\end{solution}

\begin{problem}[ODE preview]\label{prob:ii16-dossier-11}
How does Picard iteration arise for $y'=f(t,y)$?
\end{problem}
\begin{solution}
Rewrite as $y(t)=y_0+\int f(s,y(s))ds$ and iterate the integral operator.
\end{solution}

\begin{problem}[Contraction synthesis]\label{prob:ii16-dossier-12}
Why is this theorem algorithmic?
\end{problem}
\begin{solution}
Its proof constructs the solution by iteration and simultaneously provides computable geometric error bounds.
\end{solution}

\section{Exercises with complete solutions}

\begin{exercise}\label{exr:ii16-01}
Find fixed point of $T(x)=x/2+1$.
\end{exercise}
\begin{hint}
Solve algebraically. Compute a Lipschitz bound for \(T\) and verify the constant is strictly below one on the chosen domain.
\end{hint}
\begin{solution}
It is $2$.
\end{solution}

\begin{exercise}\label{exr:ii16-02}
What is the contraction factor?
\end{exercise}
\begin{hint}
Coefficient of $x$. Check that the domain is complete and that \(T\) maps it into itself before invoking the contraction theorem.
\end{hint}
\begin{solution}
$1/2$.
\end{solution}

\begin{exercise}\label{exr:ii16-03}
Why unique?
\end{exercise}
\begin{hint}
Use fixed-point inequality. Use \(d(x_n,x_*)\le q^n d(x_0,x_*)\) or the a posteriori geometric-tail estimate to quantify iteration error.
\end{hint}
\begin{solution}
Two fixed points have zero distance.
\end{solution}

\begin{exercise}\label{exr:ii16-04}
Does a contraction have to be linear?
\end{exercise}
\begin{hint}
If the global map is not contractive, restrict to a closed invariant subset where a sharper derivative or Lipschitz bound holds.
\end{hint}
\begin{solution}
No.
\end{solution}

\begin{exercise}\label{exr:ii16-05}
What property supplies the limit?
\end{exercise}
\begin{hint}
Completeness. Uniqueness follows by applying the contraction inequality to two hypothetical fixed points.
\end{hint}
\begin{solution}
Completeness of the space.
\end{solution}

\begin{exercise}\label{exr:ii16-06}
What sequence is used?
\end{exercise}
\begin{hint}
Picard iteration. To compare iterations, the smaller contraction factor gives the faster geometric asymptotic rate.
\end{hint}
\begin{solution}
$x_{n+1}=Tx_n$.
\end{solution}

\begin{exercise}\label{exr:ii16-07}
What is the qualitative error rate?
\end{exercise}
\begin{hint}
Geometric. If \(q=1\), the theorem no longer gives convergence or uniqueness; inspect a simple isometry as a counterexample.
\end{hint}
\begin{solution}
Order $q^n$.
\end{solution}

\begin{exercise}\label{exr:ii16-08}
Can $q=1$ be used?
\end{exercise}
\begin{hint}
No strict contraction. Completeness is used to ensure the Cauchy iterate sequence has a limit inside the space.
\end{hint}
\begin{solution}
Not in Banach's theorem.
\end{solution}


% BEGIN VOL02-EXPANSION II16-exercises-01
\section{Graded supplementary exercises}
These exercises extend the preserved foundational set and are graded by mathematical role.

\subsection*{Standard computations}

\begin{exercise}[Affine fixed point]\label{exr:ii16-expand-01}
Find the fixed point of \(T(x)=(x+5)/3\).
\end{exercise}
\begin{hint}
Solve \(x=T(x)\).
\end{hint}
\begin{solution}
\(x_*=5/2\).
\end{solution}

\begin{exercise}[Contraction constant]\label{exr:ii16-expand-02}
Find a contraction constant for \(T(x)=0.7x+1\) on \(\mathbb R\).
\end{exercise}
\begin{hint}
Subtract outputs.
\end{hint}
\begin{solution}
\(q=0.7\).
\end{solution}

\begin{exercise}[First Picard steps]\label{exr:ii16-expand-03}
For \(T(x)=(x+3)/4\), compute \(x_1,x_2\) from \(x_0=0\).
\end{exercise}
\begin{hint}
Apply \(T\) twice.
\end{hint}
\begin{solution}
\(x_1=3/4\), \(x_2=15/16\).
\end{solution}

\begin{exercise}[A priori error]\label{exr:ii16-expand-04}
If \(q=1/2\) and \(d(x_1,x_0)=1\), bound \(d(x_5,x_*)\).
\end{exercise}
\begin{hint}
Use \(q^n(1-q)^{-1}d(x_1,x_0)\).
\end{hint}
\begin{solution}
The bound is \(2(1/2)^5=1/16\).
\end{solution}

\begin{exercise}[A posteriori error]\label{exr:ii16-expand-05}
If \(q=1/3\) and \(d(x_n,x_{n-1})=0.006\), bound current error.
\end{exercise}
\begin{hint}
Use \(q/(1-q)\) times the last step.
\end{hint}
\begin{solution}
The error is at most \(0.003\).
\end{solution}

\subsection*{Proofs}

\begin{exercise}[Banach uniqueness]\label{exr:ii16-expand-06}
Prove a contraction has at most one fixed point.
\end{exercise}
\begin{hint}
Apply the contraction inequality to two fixed points.
\end{hint}
\begin{solution}
\(d(x,y)\le qd(x,y)\) with \(q<1\), so \(d(x,y)=0\).
\end{solution}

\begin{exercise}[Picard Cauchy]\label{exr:ii16-expand-07}
Prove Picard iterates form a Cauchy sequence.
\end{exercise}
\begin{hint}
Bound successive differences geometrically and sum a tail.
\end{hint}
\begin{solution}
\(d(x_{n+k},x_n)\le d(x_1,x_0)\sum_{j=n}^{n+k-1}q^j\), whose tail tends to zero.
\end{solution}

\begin{exercise}[Fixed point at the limit]\label{exr:ii16-expand-08}
Why is the Picard limit fixed?
\end{exercise}
\begin{hint}
Use continuity of a contraction.
\end{hint}
\begin{solution}
If \(x_n\to x_*\), then \(x_{n+1}=Tx_n\to Tx_*\), but also \(x_{n+1}\to x_*\), so \(Tx_*=x_*\).
\end{solution}

\begin{exercise}[Perturbation bound]\label{exr:ii16-expand-09}
Derive \(d(x_T,x_S)\le(1-q)^{-1}\sup d(Tx,Sx)\) for contractions with common factor \(q\).
\end{exercise}
\begin{hint}
Insert \(Tx_S\).
\end{hint}
\begin{solution}
\(d(x_T,x_S)\le qd(x_T,x_S)+\sup d(Tx,Sx)\); rearrange.
\end{solution}

\subsection*{Counterexamples and hypothesis testing}

\begin{exercise}[Completeness required]\label{exr:ii16-expand-10}
Give a contraction on an incomplete space with no fixed point.
\end{exercise}
\begin{hint}
Use \((0,1)\) and a map whose fixed point is \(1\).
\end{hint}
\begin{solution}
\(T(x)=(x+1)/2\) is a contraction on \((0,1)\) but has no fixed point there.
\end{solution}

\begin{exercise}[Strict contraction required]\label{exr:ii16-expand-11}
Give a nonexpansive self-map with no fixed point.
\end{exercise}
\begin{hint}
Use translation on \(\mathbb R\).
\end{hint}
\begin{solution}
\(T(x)=x+1\) has Lipschitz constant \(1\) and no fixed point.
\end{solution}

\begin{exercise}[Self-map required]\label{exr:ii16-expand-12}
Why can a contractive formula fail Banach's theorem if it leaves the chosen domain?
\end{exercise}
\begin{hint}
The theorem requires \(T:X\to X\).
\end{hint}
\begin{solution}
Without invariance, the iteration may exit the complete set on which the estimate was checked.
\end{solution}

\subsection*{Applications and investigations}

\begin{exercise}[Iteration stopping]\label{exr:ii16-expand-13}
How can the a posteriori bound be used as a stopping rule?
\end{exercise}
\begin{hint}
Stop once the bound computed from the latest step is below tolerance.
\end{hint}
\begin{solution}
It certifies the unknown fixed-point error using observable iterate differences.
\end{solution}

\begin{exercise}[Nonlinear equation]\label{exr:ii16-expand-14}
How does solving \(x=\cos x\) become a fixed-point problem?
\end{exercise}
\begin{hint}
Use \(T(x)=\cos x\) on a suitable invariant interval.
\end{hint}
\begin{solution}
Choose an interval where \(|\sin x|\le q<1\); Banach then proves a unique local fixed point and convergence of iteration.
\end{solution}

\subsection*{Challenge problems}

\begin{exercise}[Choose an invariant interval]\label{exr:ii16-expand-15}
Show \(T(x)=\frac12\cos x\) is a contraction mapping of \([-1,1]\) into itself.
\end{exercise}
\begin{hint}
Bound both output size and derivative.
\end{hint}
\begin{solution}
\(|T(x)|\le1/2\), so the interval is invariant, and \(|T'(x)|\le1/2\), so it is a contraction.
\end{solution}

\begin{exercise}[Optimal relaxation]\label{exr:ii16-expand-16}
For \(T_\alpha(x)=(1-\alpha)x+\alpha c\), \(0<\alpha<2\), find the contraction factor and discuss the best \(\alpha\).
\end{exercise}
\begin{hint}
Subtract outputs.
\end{hint}
\begin{solution}
The factor is \(|1-\alpha|\), minimized at \(\alpha=1\), which reaches the fixed point \(c\) in one step.
\end{solution}
% END VOL02-EXPANSION II16-exercises-01
\section*{Chapter summary}
The chapter's tools now form part of the canonical Volume II analysis toolkit and will be used in later approximation, fixed-point, and convergence arguments.
